% Generated draft: requires mathematical and source review.
\documentclass[11pt]{article}
\usepackage[margin=25mm]{geometry}
\usepackage{fontspec}
\setmainfont{Latin Modern Roman}
\usepackage{amsmath,amssymb,mathtools,unicode-math}
\setmathfont{Latin Modern Math}
\usepackage{enumitem,needspace,xurl,hyperref}
\setlength{\parindent}{0pt}
\setlength{\parskip}{5pt}
\setlength{\emergencystretch}{3em}
\newcommand{\N}{\mathbb{N}}
\newcommand{\Z}{\mathbb{Z}}
\newcommand{\Q}{\mathbb{Q}}
\newcommand{\R}{\mathbb{R}}
\newcommand{\C}{\mathbb{C}}
\newcommand{\F}{\mathbb{F}}
\newcommand{\A}{\mathbb{A}}
\newcommand{\PP}{\mathbb{P}}
\newcommand{\E}{\mathbb{E}}
\newcommand{\eps}{\varepsilon}
\newcommand{\dd}{\,\mathrm{d}}
\begin{document}
\section*{UnsolvedMath ID: 30004415}
\textbf{Source problem code:} catalogue entry
\subsection{Definitions and mathematical statement}
Let \(p\) be a prime and \(G\) a finite group. A Sylow \(p\)-subgroup \(P\leq G\) has order equal to the largest power of \(p\) dividing \(|G|\). It is abelian if \(xy=yx\) for all \(x,y\in P\). All Sylow \(p\)-subgroups of \(G\) are conjugate, so this property does not depend on the choice of \(P\). A \(p\)-element is an element of order \(p^a\) for some integer \(a\geq0\), including the identity. Write \(\operatorname{Irr}(G)\) for the irreducible complex characters of \(G\).

The supplied question asks whether the \(p\)-element submatrix \(X_p(G)\) of the character table determines whether a Sylow \(p\)-subgroup is abelian, for every finite group \(G\) and every prime \(p\). However, the supplied record does not define this submatrix or specify which labels and auxiliary data it retains, and its supplied literature assessment explicitly describes the extracted question as malformed.

A conditional interpretation is to retain every irreducible-character row and precisely the columns corresponding to conjugacy classes of \(p\)-elements. Choosing representatives \(g_1,\ldots,g_t\) of these classes would give
\[
X_p(G)=\bigl(\chi(g_j)\bigr)_{\chi\in\operatorname{Irr}(G),\,1\leq j\leq t}.
\]
If the retained data consist only of these complex entries, with row and column order disregarded, a precise candidate question is the following: for every prime \(p\), every pair of finite groups \(G,H\), and Sylow \(p\)-subgroups \(P\leq G\) and \(Q\leq H\), does
\[
X_p(G)\cong X_p(H)
\quad\Longrightarrow\quad
\bigl(P\text{ is abelian}\ \Longleftrightarrow\ Q\text{ is abelian}\bigr)
\]
hold, where \(\cong\) means equality of matrices after independent permutations of rows and columns? This is an explicitly conditional formalization; the supplied evidence does not establish that this definition of the data and their equivalence is the intended one.

The source question additionally asserts determination when \(p\notin\{3,5\}\). Its supplied dated assessment says that the cited results do not establish the claimed general determination for the extracted submatrix question. Thus the assertion for \(p\notin\{3,5\}\) is preserved as a source claim requiring clarification, rather than adopted as an established theorem. A definitive formulation, including whether the intended remaining cases are exactly \(p=3\) and \(p=5\), requires human review of the source.
\subsection{Short English statement}
The supplied question asks whether character-table data on elements of prime-power order detect an abelian Sylow subgroup. A conditional matrix formulation is given above, but the source does not specify the retained data precisely. Its claim for \(p\notin\{3,5\}\), and its supplied open status, require source review.
\subsection{Sources}
\begin{itemize}
\item Local UnsolvedMath registry, supplied complete source record, permanent ID \(30004415\), title ``Detecting Abelian Sylow Subgroups from Character Tables,'' source problem code \texttt{catalogue entry}; collection: Oberwolfach Reports Open Problems. The record supplies proposed year 2020 and no named proposer.
\item Supplied report attribution: \emph{Mini-Workshop: Kronecker, Plethysm, and Sylow Branching Coefficients and their Applications to Complexity Theory} (2021), Oberwolfach Reports, pp.\ 517--568, DOI 10.4171/owr/2020/9. \url{https://doi.org/10.4171/owr/2020/9}. The literature-source list labels this same link ``Oberwolfach Report 9/2020.'' The differing year indications are retained without reconciliation.
\item Supplied literature reference: Kessar, Malle, and Navarro, ``abelian Sylow subgroups.'' \url{https://arxiv.org/abs/1407.4495}.
\item Supplied literature reference: Kessar, Malle, and Navarro, ``published version.'' \url{https://doi.org/10.1016/j.jalgebra.2014.08.012}.
\item Supplied literature assessment checked on \texttt{2026-08-22}: status ``open.'' This is supplied dated curation, not independently verified. The assessment states that the full character table detects abelian Sylow subgroups, but that the extracted \(p\)-element-submatrix question is malformed and its claimed general determination is not established by the cited results. The background describes this as dated literature triage derived from the supplied Oberwolfach curation and requests independent verification of time-sensitive status and cited claims.
\end{itemize}
\end{document}
